% !TeX root = ../main.tex
% Fonte: Lezione 7.pdf, pp. 19--21.
\section{Insiemi aperti e insiemi chiusi}

\begin{definizione}[Insieme aperto]
Un insieme $A\subseteq\R$ si dice \emph{aperto} se è un'unione di intervalli aperti: esiste una famiglia $(I_\alpha)_{\alpha\in J}$ di intervalli aperti tale che
\[
A=\bigcup_{\alpha\in J}I_\alpha.
\]
\end{definizione}

\textbf{Famiglie ammesse.} La famiglia può avere un numero finito o infinito di elementi; si ammette anche la famiglia vuota, la cui unione è $\varnothing$.

\textbf{Esempi di aperti.}
\begin{itemize}
  \item $\R=(-\infty,+\infty)$ è aperto;
  \item $\varnothing$ è aperto, come unione della famiglia vuota;
  \item $\R\setminus\Z=\bigcup_{n\in\Z}(n,n+1)$ è aperto;
  \item $\R^*=\R\setminus\{0\}=(-\infty,0)\cup(0,+\infty)$ è aperto.
\end{itemize}
L'insieme di tutti gli aperti si chiama \emph{topologia usuale} di $\R$.

\begin{definizione}[Insieme chiuso]
Un insieme $B\subseteq\R$ si dice \emph{chiuso} se il complemento $B^c=\R\setminus B$ è aperto.
\end{definizione}

\textbf{Esempi di chiusi.} Per verificare che un insieme è chiuso basta quindi descrivere il complemento come unione di intervalli aperti:
\begin{itemize}
  \item $\R$ è chiuso, perché $\R^c=\varnothing$ è aperto;
  \item $\varnothing$ è chiuso, perché $\varnothing^c=\R$ è aperto;
  \item per $a\le b$, $[a,b]$ è chiuso, perché $[a,b]^c=(-\infty,a)\cup(b,+\infty)$;
  \item $\{0\}$ è chiuso, perché il suo complemento è $(-\infty,0)\cup(0,+\infty)$;
  \item $[a,+\infty)$ è chiuso, perché ha complemento $(-\infty,a)$.
\end{itemize}
\textbf{Rapporto fra le due proprietà.} Aperto e chiuso non sono proprietà mutuamente esclusive: $\varnothing$ e $\R$ le possiedono entrambe. Esistono anche insiemi che non possiedono nessuna delle due, come vedremo per $[0,1)$.

\section{Il criterio locale di apertura}

\begin{proposizione}[Caratterizzazione degli aperti]
\label{prop:aperti-criterio}
Un insieme $A\subseteq\R$ è aperto se e solo se
\[
\forall x\in A\ \exists\varepsilon>0:
\quad (x-\varepsilon,x+\varepsilon)\subseteq A.
\]
\end{proposizione}
\begin{proof}
\passomath{\Rightarrow} Se $A$ è un'unione di intervalli aperti e $x\in A$, allora $x$ appartiene a uno di tali intervalli, diciamo $I$. Poiché $x$ non è un estremo di $I$, si può scegliere $\varepsilon>0$ abbastanza piccolo affinché $(x-\varepsilon,x+\varepsilon)\subseteq I\subseteq A$. Per $I=(a,b)$ con estremi finiti basta, per esempio, $\varepsilon=\min\{x-a,b-x\}/2$; per una semiretta basta considerare la distanza dall'unico estremo finito, e per $I=\R$ ogni raggio va bene.

\passomath{\Leftarrow} Per ogni $x\in A$ scegliamo un intervallo $I_x=(x-\varepsilon_x,x+\varepsilon_x)\subseteq A$ come nell'ipotesi. Allora
\[
A=\bigcup_{x\in A}I_x:
\]
\passomath{A\subseteq\bigcup_{x\in A}I_x} Ogni $x\in A$ appartiene al proprio intervallo $I_x$.

\passomath{\bigcup_{x\in A}I_x\subseteq A} Tutti gli intervalli $I_x$ sono contenuti in $A$. È dunque una rappresentazione di $A$ come unione di intervalli aperti. Se $A=\varnothing$, entrambi i criteri valgono con la convenzione della famiglia vuota.
\end{proof}

La proprietà è locale: intorno a ogni punto di $A$ si deve poter ritagliare un piccolo intervallo che resti interamente in $A$. Il raggio può dipendere dal punto; non serve un unico raggio valido per tutti.

\begin{esempio}
\textbf{Non apertura.} L'insieme $A=[0,1)$ non è aperto: $0\in A$, ma ogni intervallo $(-\varepsilon,\varepsilon)$ contiene punti negativi, esterni ad $A$.

\textbf{Non chiusura.} Non è neppure chiuso, perché
\[
A^c=(-\infty,0)\cup[1,+\infty)
\]
non è aperto: contiene $1$, ma nessun intervallo aperto centrato in $1$ è contenuto nel complemento. Basta infatti scegliere, per ogni $\varepsilon>0$, il punto $1-\min\{\varepsilon/2,1/2\}\in A\cap(1-\varepsilon,1+\varepsilon)$.
\end{esempio}
